% Generated draft: requires mathematical and source review.
\documentclass[11pt]{article}
\usepackage[margin=25mm]{geometry}
\usepackage{fontspec}
\setmainfont{Latin Modern Roman}
\usepackage{amsmath,amssymb,mathtools,unicode-math}
\setmathfont{Latin Modern Math}
\usepackage{enumitem,needspace,xurl,hyperref}
\setlength{\parindent}{0pt}
\setlength{\parskip}{5pt}
\setlength{\emergencystretch}{3em}
\newcommand{\N}{\mathbb{N}}
\newcommand{\Z}{\mathbb{Z}}
\newcommand{\Q}{\mathbb{Q}}
\newcommand{\R}{\mathbb{R}}
\newcommand{\C}{\mathbb{C}}
\newcommand{\F}{\mathbb{F}}
\newcommand{\A}{\mathbb{A}}
\newcommand{\PP}{\mathbb{P}}
\newcommand{\E}{\mathbb{E}}
\newcommand{\eps}{\varepsilon}
\newcommand{\dd}{\,\mathrm{d}}
\begin{document}
\section*{UnsolvedMath ID: 30006520}
\textbf{Source problem code:} catalogue entry
\subsection{Definitions and mathematical statement}
The supplied record asks the following support question. It does not explicitly specify the version of Grothendieck polynomials or the ambient symmetric group. The formulation below uses the ordinary single type~A Grothendieck polynomials and the stable symmetric group to make the notation precise; these conventions require confirmation against the source.

Let \(S_\infty=\bigcup_{n\geq 1}S_n\), where the inclusions fix all additional letters. For \(u\in S_\infty\), let
\[
\ell(u)=\#\{(i,j):i<j,\ u(i)>u(j)\}
\]
be its Coxeter length. Define the left weak order by
\[
u\leq_L v
\quad\Longleftrightarrow\quad
\text{there exists }a\in S_\infty\text{ with }v=au
\text{ and }\ell(v)=\ell(a)+\ell(u),
\]
and write \(u<_L v\) when additionally \(u\ne v\).

Let \(\mathfrak S_u\) denote the ordinary type~A Schubert polynomial indexed by \(u\). In the ordinary single Grothendieck polynomial convention, expand \(G_w\) in the Schubert basis of \(\Z[x_1,x_2,\ldots]\):
\[
G_w=\sum_{u\in S_\infty}c_{w,u}\mathfrak S_u,
\qquad c_{w,u}\in\Z.
\]
The sum has finite support. Define its Schubert support and total polynomial degree by
\[
\operatorname{Supp}_S(G_w)=\{u\in S_\infty:c_{w,u}\ne0\},
\qquad
\deg(G_w)=\max\{|\alpha|:[x^\alpha]G_w\ne0\},
\]
where \(\alpha\) ranges over finitely supported sequences of nonnegative integers, \(x^\alpha=\prod_i x_i^{\alpha_i}\), and \(|\alpha|=\sum_i\alpha_i\).

\textbf{Question (supplied Conjecture 1).} With the intended conventions fixed, is it true that for every indexing permutation \(w\) and every \(u\in\operatorname{Supp}_S(G_w)\),
\[
\ell(u)<\deg(G_w)
\quad\Longrightarrow\quad
\exists v\in\operatorname{Supp}_S(G_w)\text{ such that }u<_L v?
\]
Here the displayed stable interpretation takes both indexing permutations in \(S_\infty\). The supplied record does not resolve whether a fixed finite symmetric group was intended instead. The conclusion requires some strictly larger element in left weak order; it does not state that every larger permutation belongs to the support, nor does it require that \(v\) cover \(u\).

The excerpt also mentions an unspecified stronger statement and analogues of Grothendieck monomial support conjectures cited there as reference~[3]. Neither that stronger statement nor the identity of reference~[3] is supplied, so neither is incorporated into this formulation.
\subsection{Short English statement}
The question asks whether every permutation in the Schubert expansion of a Grothendieck polynomial whose Coxeter length is below the polynomial's total degree has a strictly larger permutation, in left weak order, in the same support. The polynomial convention and ambient symmetric group require source confirmation. The record labels the problem partially solved; that label is supplied dated curation, not an independently verified mathematical assessment.
\subsection{Sources}
\begin{itemize}
\item Local UnsolvedMath registry, \texttt{data/unsolvedmath/problems.json}, supplied record with permanent ID \(30006520\), title ``Upward Closure of Grothendieck Schubert Supports,'' exact source problem code \texttt{catalogue entry}. The record identifies the excerpt as Conjecture~1, gives proposed year \(2026\), and supplies no individual proposer.
\item Supplied report citation: \emph{Enumerative Combinatorics} (2026), Oberwolfach Reports, pp.~81--156, DOI \texttt{10.4171/owr/2026/2}. \url{https://doi.org/10.4171/owr/2026/2}. The background describes Oberwolfach Reports as documenting mathematical workshops and their open-problem discussions. The supplied citation covers the report; no precise page or contribution author for this conjecture is supplied.
\item Mészáros--Setiabrata--St. Dizier, \emph{On the support of Grothendieck polynomials}. Supplied literature source: \url{https://arxiv.org/abs/2201.09452}.
\item Chou--Yu, \emph{Grothendieck polynomials of inverse fireworks permutations}. Supplied literature source: \url{https://arxiv.org/abs/2403.12186}.
\item Supplied literature assessment, checked \texttt{2026-08-22}: status ``partially solved''; the general Grothendieck-support assertion is described as remaining open, with special cases including Grassmannian and fireworks families proved. This is supplied dated curation derived from the OWR record, not independently verified. The supplied background expressly calls for independent verification of time-sensitive status and cited claims.
\item The registry labels the statement \texttt{corrected\_verified}, but its accompanying qualification says only formula delimiters, notation, grammar, and typography were repaired, with mathematical content and scope preserved. This label does not establish independent mathematical verification.
\end{itemize}
\end{document}
